% Table 1. Descriptive, not empirical: every cell is a provision read off the
% primary documents cited in the note, so there is no frozen output behind it
% and make_tables.py does not generate it. Values must be changed only against
% the source documents.
\begin{table}[tbp]
\caption{The two published rules, side by side.}
\label{tab:rules}
\centering\small
\setlength{\tabcolsep}{5pt}
\begin{tabularx}{\textwidth}{@{}lYY@{}}
\toprule
 & Optimised index & Heuristic index\\
\midrule
Objective        & Minimise predicted variance & None stated\\
Risk input       & Full covariance matrix from a commercial factor model
                 & Trailing one-year realised volatility, per stock\\
Selection        & Implied by the programme & The 100 lowest-volatility
                   parent members\\
Weighting        & Implied by the programme & Proportional to
                   $1/\sigma_i$\\
\addlinespace
Weight cap       & $\min(1.5\%,\ 20\times$ parent weight$)$ & None\\
Minimum holding  & 0.05\%, or zero & None\\
Sector band      & $\pm5\%$ active per sector & None\\
Style band       & $\pm0.25$ SD active per style, volatility exempt & None\\
Turnover cap     & 10\% one-way per review & None\\
\addlinespace
Review frequency & Semi-annual & Quarterly\\
Effective date   & Last business day of May and November
                 & Third Friday of the review month\\
Risk-model date  & End of the month preceding the review & Not applicable\\
Infeasibility    & Documented relaxation sequence & Cannot arise\\
\bottomrule
\end{tabularx}
\notetext{Provisions for the optimised index are from its provider's published
methodology \citep{msci2012minvol}; those for the heuristic index are from the
tracking fund's summary prospectus \citep{invesco2024splv497k}, which states the
selection and weighting rules in the issuer's own words. The heuristic index's
review calendar and effective date are the one entry in this table taken from
published convention rather than from a methodology document: the provider's
methodology PDF refused the request from the run environment and was not
retrieved, which Section~\ref{sec:limits} records as an unverified parameter.
The style band is the row that does not survive translation into another risk
model, and Section~\ref{sec:rules} takes it up.}
\end{table}
